\begin{frame}{Index 1 of 4: principle, slide, prompt}
\scriptsize
\begin{tabular}{@{}r>{\raggedright\arraybackslash}p{5.4cm}>{\raggedright\arraybackslash}p{4.6cm}p{1.2cm}@{}}
\textbf{Slide} & \textbf{Title} & \textbf{Principle} & \textbf{Prompt}\\\midrule
5 & Union to library & ASYMMETRY & \\
6 & Give it a name: a directed, weighted graph & MODEL VS PROCEDURE & \\
7 & A six-sided die gives six cities & ASYMMETRY & \\
9 & A tour: visit every city once, then come home & MODEL VS PROCEDURE & \\
10 & Total the six legs, step by step & CHECK THE WORK & \\
11 & What are we optimizing? & MODEL VS PROCEDURE & \\
12 & Idea 1: greedy, always go to the nearest city & GREEDY IS NOT OPTIMAL & \\
13 & Greedy, worked on the table from city 1 & GREEDY IS NOT OPTIMAL & \\
14 & Where greedy is fooled & GREEDY IS NOT OPTIMAL & \\
\end{tabular}
\note{Use this to find a slide fast: slide number, which lesson idea it shows, and which prompt file it pairs with (prompts/NN, a = first block, b = second block).}
\end{frame}
\begin{frame}{Index 2 of 4: principle, slide, prompt}
\scriptsize
\begin{tabular}{@{}r>{\raggedright\arraybackslash}p{5.4cm}>{\raggedright\arraybackslash}p{4.6cm}p{1.2cm}@{}}
\textbf{Slide} & \textbf{Title} & \textbf{Principle} & \textbf{Prompt}\\\midrule
15 & Idea 2: guess and check with the die & CHECK THE WORK & \\
16 & What rejecting repeats costs & COUNTING / PATTERN & \\
17 & A short worked trace & CHECK THE WORK & \\
18 & Guessing can find the best, but cannot tell you it did & CHECK THE WORK & \\
19 & Count the tours: 2 cities, then 3 & COUNTING / PATTERN & \\
20 & 4 cities & COUNTING / PATTERN & \\
21 & The reasoning: choices shrink by one each step & COUNTING / PATTERN & \\
22 & The pattern & COUNTING / PATTERN & \\
23 & Which counting convention? Say it once & COUNTING / PATTERN & \\
\end{tabular}
\note{Use this to find a slide fast: slide number, which lesson idea it shows, and which prompt file it pairs with (prompts/NN, a = first block, b = second block).}
\end{frame}
\begin{frame}{Index 3 of 4: principle, slide, prompt}
\scriptsize
\begin{tabular}{@{}r>{\raggedright\arraybackslash}p{5.4cm}>{\raggedright\arraybackslash}p{4.6cm}p{1.2cm}@{}}
\textbf{Slide} & \textbf{Title} & \textbf{Principle} & \textbf{Prompt}\\\midrule
24 & Plot it (linear) & COUNTING / PATTERN & \\
25 & Plot it (log) & COUNTING / PATTERN & \\
26 & What the curve means for brute force & COUNTING / PATTERN & \\
27 & Then we started writing code & MODEL VS PROCEDURE & \\
28 & Step 1: read the table & MODEL VS PROCEDURE & 01a\\
29 & Step 2: the length of a tour & CHECK THE WORK & 02a\\
30 & Step 3: check 1-2-3-4-5-6-1, then write tests & CHECK THE WORK & 03a\\
31 & Step 3b: tests that guard the model & ASYMMETRY & 03b\\
32 & Step 4: greedy & GREEDY IS NOT OPTIMAL & 04a\\
\end{tabular}
\note{Use this to find a slide fast: slide number, which lesson idea it shows, and which prompt file it pairs with (prompts/NN, a = first block, b = second block).}
\end{frame}
\begin{frame}{Index 4 of 4: principle, slide, prompt}
\scriptsize
\begin{tabular}{@{}r>{\raggedright\arraybackslash}p{5.4cm}>{\raggedright\arraybackslash}p{4.6cm}p{1.2cm}@{}}
\textbf{Slide} & \textbf{Title} & \textbf{Principle} & \textbf{Prompt}\\\midrule
33 & Step 5: guess and check with random & CHECK THE WORK & 05a\\
34 & Step 6: brute force over permutations & COUNTING / PATTERN & 06a\\
35 & Step 6b: assert the count matches the formula & COUNTING / PATTERN & 06b\\
36 & Model versus procedure: each step, again & MODEL VS PROCEDURE & \\
37 & If the street were two-way & MODEL VS PROCEDURE & \\
39 & Gate preview & CHECK THE WORK & \\
40 & Exit trace & CHECK THE WORK & \\
\end{tabular}
\note{Use this to find a slide fast: slide number, which lesson idea it shows, and which prompt file it pairs with (prompts/NN, a = first block, b = second block).}
\end{frame}
